\documentclass[11pt,a4paper]{article}
\usepackage[utf8]{inputenc}
\usepackage[T1]{fontenc}
\usepackage[brazil]{babel}
\usepackage[a4paper,margin=2.3cm]{geometry}
\usepackage{enumitem}
\usepackage{hyperref}
\usepackage{xcolor}
\setlength{\parindent}{0pt}
\setlength{\parskip}{5pt}
\sloppy
\setlength{\emergencystretch}{3em}
\hypersetup{colorlinks=true,urlcolor=blue,linkcolor=black}
\newcommand{\cabecalho}[2]{\begin{center}{\small \textbf{Sprint 2 -- #1}}\\[1mm]\textit{Inventário de Ativos e Vulnerabilidades}\\Matheus Moreli Durão Silva -- 12621CBS217\\Cibersegurança -- Bacharelado -- Noturno\end{center}\hrule\vspace{4mm}}
\newcommand{\bloco}[1]{\par\vspace{2mm}\textbf{#1}\par}
\newcommand{\codigo}[1]{\texttt{#1}}

\begin{document}\cabecalho{S2\_08 -- Cadastro de Ativos}{}
\bloco{Dados mínimos} O cadastro pede ID único, nome/hostname, responsável, localização, tipo e criticidade. Não pede senha, token ou outra credencial: inventário identifica o ativo, não armazena segredos.
\bloco{Lógica} A interface recebe as respostas. Depois \codigo{cadastrar\_ativo} verifica o ID, chama as validações de texto e criticidade, cria um objeto \codigo{Ativo}, registra ``Ativo cadastrado'' no histórico e só então o insere no dicionário. Se ocorrer erro antes dessa linha, não há ativo incompleto na base.
\bloco{Criticidade} As opções padronizadas são Baixa, Média, Alta e Crítica. Essa classificação é do ativo; ela ajuda a priorizar uma vulnerabilidade com base no contexto, e não apenas na nota técnica CVSS.
\bloco{Como testar} Cadastre ID 10 com criticidade Alta e busque pelo ID. Tente cadastrar o mesmo ID: deve ser recusado. Consulte o histórico do registro criado.
\bloco{Checklist atendido} ID único, obrigatórios validados, criticidade padronizada, nenhuma credencial solicitada e ativo consultável após criação.
\end{document}
